\fiche{25}{Se repérer : scripts, imports et notebook}{\src{02_dataset.py}{load\_local\_module}}
\objectif{Identifier le chemin réellement exécuté avant d'expliquer les données.}
Le dépôt propose une progression lisible : tokenisation, jeux de données,
attention, blocs Transformer, modèle, entraînement, génération. Cette progression
est pédagogique, mais elle n'impose pas que chaque programme utilise exactement
la démonstration du précédent. Pour suivre une expérience, partons de son point
d'entrée, puis descendons vers les fonctions appelées. Le répertoire de référence
est \path{/home/runner/work/SHOKOBO/SHOKOBO/mini_gpt}.

\notion{Un fichier Python n'est pas toujours importable par une instruction ordinaire.}
Les noms numérotés commencent par des chiffres : écrire une instruction
d'importation avec ce nom comme identifiant Python serait invalide.
La fonction étudiée construit donc un chemin depuis \code{BASE\_DIR},
obtient une spécification avec \code{importlib.util.spec\_from\_file\_location},
crée un objet module, puis exécute son contenu avec
\code{spec.loader.exec\_module}. Le nom logique donné au module ne doit pas
être confondu avec son nom de fichier. Ce mécanisme charge les définitions,
et exécute aussi les instructions placées au niveau supérieur.

La garde \code{if \_\_name\_\_ == '\_\_main\_\_'} empêche en revanche de
lancer automatiquement la démonstration lorsque le fichier est chargé comme
module. Ainsi, récupérer \code{NextTokenDataset} ne signifie pas afficher les
fenêtres de son \code{main}. Le chemin est calculé relativement à
\code{\_\_file\_\_}, non relativement au dossier depuis lequel on lance Python.

\notion{Deux parcours de données coexistent.}
La démonstration de \src{02_dataset.py}{main} lit tout le petit texte,
apprend un tokenizer caractère et emploie \code{NextTokenDataset}.
L'entraînement actif de \src{06_train.py}{train} appelle
\code{prepare\_corpus}, puis \code{MemmapTokenDataset}. La fonction historique
\code{create\_dataloaders} existe encore, mais ce n'est pas le chemin appelé
par \code{train}. Lire seulement son découpage à 90\,\% donnerait donc une
description incomplète du pipeline actuel.

Le notebook
\path{/home/runner/work/SHOKOBO/SHOKOBO/mini_gpt/mini_gpt_pedagogique.ipynb}
sert à explorer interactivement. L'ordre d'exécution des cellules peut laisser
des objets anciens en mémoire ; un script relancé repart dans un nouveau
processus. Le guide
\path{/home/runner/work/SHOKOBO/SHOKOBO/mini_gpt/README.md}
décrit ces deux usages sans remplacer la lecture des appels effectifs.

\exercice{Un lecteur affirme : « L'entraînement mélange une liste de tous les IDs
avec la fonction historique. » Comment vérifier cette affirmation ?}
\corrige{Suivre \code{main}, puis \code{train} dans le script d'entraînement.
On rencontre la préparation des fichiers, les datasets memmap et un
\code{RandomSampler} avec remise. La fonction historique n'est pas appelée.
Il faut distinguer une définition présente dans le fichier d'une fonction
effectivement exécutée : la proximité visuelle n'établit aucune relation d'appel.}

\fiche{26}{Texte Unicode, octets et vocabulaire}{\src{01_tokenizer.py}{train\_tokenizer}}
\objectif{Distinguer ce que l'on voit, ce que Python compte et ce que le disque stocke.}
Un texte Python de type \code{str} est une suite de points de code Unicode.
Un fichier UTF-8 est une suite d'octets. Un octet possède 256 valeurs possibles ;
un point de code peut demander plusieurs octets lors de son encodage UTF-8.
Ni l'un ni l'autre ne correspond toujours à une lettre visuelle complète :
un accent combinant ou plusieurs signes composant un emoji peuvent participer
à un seul graphème perçu par le lecteur.

\notion{La longueur dépend de l'unité choisie.}
Le texte original « Aé🙂 » contient trois points de code. En UTF-8, leurs
tailles respectives sont un, deux et quatre octets. On obtient donc :
\[
    N_{\mathrm{caractères}}=3,\qquad N_{\mathrm{octets}}=1+2+4=7.
\]
Le tokenizer caractère produit ici trois IDs, mais leur stockage ultérieur
en entiers de huit octets occupera vingt-quatre octets, hors métadonnées.
Transformer du texte en nombres n'est donc pas automatiquement le compresser.
La tokenisation prépare une interface pour le modèle ; sa taille disque
dépend également du type numérique choisi.

\notion{Aucune normalisation Unicode n'est appliquée.}
Un « é » précomposé est un seul point de code ; un « e » suivi d'un accent
aigu combinant en comporte deux. Ces formes peuvent sembler identiques à
l'écran, mais \code{set(text)} et \code{encode} les traitent différemment.
Le code ne les convertit pas vers une forme commune, ne met pas le texte en
minuscules et ne retire pas les espaces. Une politique de nettoyage serait
une décision supplémentaire, absente de cette implémentation.

Le vocabulaire est une association entre unités et entiers, pas un dictionnaire
de définitions linguistiques. Sa taille est notée \(V\).
Un ID valide appartient à \(\{0,\ldots,V-1\}\).
L'ID d'une lettre n'exprime ni sa fréquence ni sa proximité avec une autre
lettre : ces entiers désigneront des lignes d'une table d'embeddings.
La dimension \(D\) de cette table sera une autre quantité, étudiée ensuite.

\attention{Dans ce chapitre, « caractère » désigne le point de code manipulé
par Python, sauf mention explicite d'un graphème.}
Cette convention rend les calculs reproductibles, notamment pour le découpage
du corpus. Compter les octets du fichier ne suffit pas à retrouver la frontière
de validation calculée en caractères.

\exercice{Comparer « é » et la séquence \code{e} suivie de U+0301 :
combien de caractères et d'octets UTF-8 chacune contient-elle ?}
\corrige{La première contient un caractère et deux octets. La seconde contient
deux caractères et trois octets : un pour \code{e}, deux pour l'accent.
Sans normalisation, le tokenizer caractère peut donc produire une ou deux
positions pour deux affichages visuellement voisins.}

\fiche{27}{Construire le tokenizer caractère}{\src{01_tokenizer.py}{SimpleTokenizer}}
\objectif{Reconstruire exactement le vocabulaire, y compris l'entrée inconnue.}
Le constructeur accepte un itérable de chaînes. Dans l'usage caractère,
cet itérable fournit les caractères du texte ou l'alphabet déjà collecté.
La première opération est \code{sorted(set(vocab) - \{'<unk>'\})}.
L'ensemble élimine les répétitions ; le tri fixe un ordre indépendant de
l'ordre de découverte. Pour des caractères isolés, ce tri suit l'ordre Unicode,
et non l'ordre alphabétique d'un dictionnaire français.

\notion{Deux tables réciproques rendent les conversions explicites.}
La liste \code{itos} associe un indice à un token. Le dictionnaire
\code{stoi} associe un token à son indice. Avant les caractères triés,
le constructeur ajoute l'entrée spéciale \code{<unk>}. Elle reçoit donc
toujours l'ID zéro, enregistré dans \code{unk\_id}.
Cette entrée représente une information inconnue, pas une fin de phrase,
un remplissage de batch ou un début de document.

Prenons le texte original « baba é ». Son alphabet comporte une espace,
\code{a}, \code{b} et \code{é}. La liste construite est :
\[
    \mathrm{itos}=[\texttt{<unk>},\ \text{espace},\ a,\ b,\ \text{é}],
    \qquad V=5.
\]
Le \code{b} reçoit l'ID 3 même s'il apparaît en premier.
Le mot « baba » s'encode donc en \(3,2,3,2\).
Répéter ce texte mille fois ne change pas le vocabulaire :
les fréquences seront utiles au modèle, mais elles n'interviennent pas
dans cette construction particulière.

\notion{Apprentissage par flux et apprentissage depuis une chaîne.}
\code{from\_text} transmet directement le texte au constructeur.
\code{train\_tokenizer}, en mode \code{char}, parcourt au contraire les blocs
successifs et enrichit un ensemble \code{alphabet}. Il refuse un alphabet vide,
puis construit le même type de tokenizer. L'ordre des blocs ne change pas les
IDs si leur ensemble de caractères reste identique. Le paramètre
\code{vocab\_size} est ignoré dans cette branche : on ne tronque pas l'alphabet
à une taille demandée.

\attention{Le marqueur spécial est une chaîne entière.}
Lorsque le corpus contient littéralement les cinq caractères de
\code{<unk>}, le mode caractère les parcourt séparément. Ils ne sont pas
reconnus comme une occurrence magique du token spécial.
Le retrait dans le constructeur évite surtout de dupliquer ce marqueur
si un vocabulaire fourni le contient déjà comme entrée.

\exercice{Construire le vocabulaire de « cabac », puis encoder « bac ».
Que change une demande de vocabulaire de taille 2 en mode caractère ?}
\corrige{On obtient \code{['<unk>', 'a', 'b', 'c']}, donc \(V=4\), et
« bac » devient \([2,1,3]\). La taille demandée ne change rien :
les trois caractères observés et l'entrée inconnue sont conservés.}

\fiche{28}{Encoder, décoder et perdre une information}{\src{01_tokenizer.py}{SimpleTokenizer}}
\objectif{Savoir quand l'aller-retour restitue le texte et quand il ne le peut pas.}
Encoder consiste à parcourir le texte caractère par caractère. Pour chacun,
\code{stoi.get(char, unk\_id)} fournit soit son ID connu, soit zéro.
Décoder parcourt les IDs, cherche les tokens correspondants, puis concatène
leurs représentations. Dans cette dernière opération, l'entrée spéciale
\code{<unk>} est rendue par un point d'interrogation.

\notion{L'aller-retour est exact sur l'alphabet connu.}
Construisons le tokenizer avec le texte original « ab a ».
Sa table est \code{['<unk>', ' ', 'a', 'b']}.
L'encodage de « ba a » donne \([3,2,1,2]\), dont le décodage redonne
exactement « ba a ». L'espace fait partie des données au même titre que les
lettres. Il n'est ni supprimé ni reconstruit par une règle grammaticale.
\[
    \operatorname{decode}(\operatorname{encode}(s))=s
    \quad\text{si tous les caractères de }s\text{ sont connus.}
\]

En revanche, « bz » devient \([3,0]\), puis « b? ».
Une chaîne « bΩ » suit le même trajet. Deux entrées distinctes ayant la même
image, aucune fonction de décodage ne peut déterminer quel caractère inconnu
était présent. Cette perte n'est pas une erreur du réseau : elle se produit
avant son intervention. Elle illustre une application non injective,
notion mathématique utile pour comprendre les limites d'une représentation.

\notion{Les méthodes d'inspection ne rendent pas toutes la même chaîne.}
\code{id\_to\_token(0)} retourne le marqueur \code{<unk>},
tandis que \code{decode([0])} retourne \code{?}.
De même, un entier négatif ou supérieur au dernier indice valide est converti
en marqueur inconnu lors du décodage. Le test de borne empêche notamment
qu'un indice négatif sélectionne le dernier élément, comme le ferait
habituellement une liste Python.

\attention{Un vrai point d'interrogation peut lui-même appartenir au vocabulaire.}
L'affichage décodé ne permet alors pas de distinguer ce caractère d'un inconnu
remplacé. Pour mesurer les inconnus, il faut compter \code{unk\_id} dans les
IDs, pas compter les points d'interrogation du texte reconstruit.
Le taux obtenu décrit la couverture de l'alphabet d'entraînement.
Sur une validation contenant une nouvelle écriture, il peut être élevé,
même si le programme fonctionne exactement comme prévu.

\exercice{Avec le vocabulaire précédent, décoder \([2,-1,99,3]\), puis
donner le résultat de \code{id\_to\_token(-1)}. Peut-on retrouver les inconnus ?}
\corrige{Le texte est « a??b ». La méthode d'inspection retourne
\code{<unk>}. Les deux IDs invalides ne révèlent aucun caractère initial ;
le décodage fournit une représentation de remplacement, pas une réparation.
Il faut conserver le texte brut si cette information doit rester accessible.}

\fiche{29}{Sauvegarder l'état caractère sans réapprendre}{\src{01_tokenizer.py}{tokenizer\_from\_state}}
\objectif{Comprendre pourquoi les poids du modèle ne suffisent pas à interpréter ses IDs.}
Une ligne d'embedding correspond à un ID précis. Si l'on recharge les poids
avec un vocabulaire différent, les mêmes nombres peuvent désigner d'autres
caractères. Un modèle ayant associé la ligne 2 à \code{a} ne reconnaît pas
spontanément que cette ligne désigne désormais \code{b}.
Sauvegarder l'état du tokenizer fait donc partie de la sauvegarde du sens
des données, et pas seulement du confort de programmation.

\notion{Un état compact suffit pour le mode caractère.}
\code{to\_state} renvoie un dictionnaire comportant \code{kind: 'char'} et
la liste \code{vocab}, égale à \code{itos[1:]}.
Le marqueur inconnu est omis parce que le constructeur le réintroduira en
position zéro. La restauration appelle \code{SimpleTokenizer(state['vocab'])},
sans relire le corpus ni réapprendre quoi que ce soit.
Pour un état produit par ce constructeur, le tri reconstruit exactement
les mêmes tables.

\begin{lstlisting}[language=Python]
state = tokenizer.to_state()
restored = tokenizer_from_state(state)
assert restored.encode("aba") == tokenizer.encode("aba")
\end{lstlisting}

Ce court exemple suppose les objets déjà chargés. L'égalité concerne des
IDs, ce qui est plus exigeant que la seule égalité de taille du vocabulaire.
Deux vocabulaires peuvent avoir dix entrées et pourtant attribuer des
significations différentes à chacun des dix indices.

Dans \src{02_dataset.py}{prepare\_corpus}, l'état est écrit en JSON UTF-8
avec \code{ensure\_ascii=False}. Les accents peuvent ainsi rester lisibles
dans le fichier
\path{/home/runner/work/SHOKOBO/SHOKOBO/mini_gpt/checkpoints/tokens/tokenizer.json}
pour la sortie par défaut. Une sérialisation JSON avec échappements Unicode
aurait le même sens après lecture ; ce choix concerne surtout la lisibilité.
Dans \src{06_train.py}{train}, le checkpoint conserve également
\code{tokenizer.to\_state()}, ainsi qu'un champ \code{vocab} historique
pour le mode caractère.

\attention{Un état enregistré n'est pas une validation complète de son contenu.}
Le restaurateur contrôle le type annoncé par \code{kind}, mais ne transforme
pas un dictionnaire arbitraire en état fiable. Il faut garder ensemble
le bon tokenizer, les fichiers d'IDs et les poids correspondants.
Modifier manuellement le vocabulaire après préparation peut rendre les fichiers
sémantiquement incohérents sans changer leur longueur en octets.

\exercice{Un état possède \code{vocab=['a', 'c']}. Quelle table sera restaurée ?
Que se passe-t-il si l'on ajoute \code{'b'} avant de recharger les mêmes IDs ?}
\corrige{La table initiale est \code{['<unk>', 'a', 'c']}.
Après ajout et tri, elle devient \code{['<unk>', 'a', 'b', 'c']}.
L'ancien ID 2, qui signifiait \code{c}, signifie alors \code{b}.
L'absence d'erreur informatique ne garantit donc pas la cohérence du modèle.}

\fiche{30}{BPE : apprendre des fusions, à la main}{\src{01_tokenizer.py}{train\_tokenizer}}
\objectif{Comprendre le principe des paires fréquentes avant les détails de la bibliothèque.}
Le mode caractère impose une position par caractère. BPE cherche à représenter
des suites fréquentes par des unités plus longues. Son nom, \emph{Byte Pair
Encoding}, rappelle une méthode fondée sur le regroupement de paires.
Pour isoler cette idée, travaillons d'abord sur une version jouet dont les
symboles initiaux sont des lettres. La version byte-level du dépôt sera
précisée à la fiche suivante.

\notion{Une fusion remplace deux symboles adjacents par un nouveau symbole.}
Considérons trois occurrences séparées du mot inventé « baba ».
Chaque occurrence commence comme \([b,a,b,a]\).
Au total, la paire \((b,a)\) apparaît six fois et la paire \((a,b)\)
trois fois. On choisit donc \((b,a)\), puis chaque occurrence devient
\([ba,ba]\). La paire \((ba,ba)\) apparaît maintenant trois fois :
une seconde fusion produit le symbole \([baba]\).
\[
    [b,a,b,a]\ \longrightarrow\ [ba,ba]\ \longrightarrow\ [baba].
\]
Les frontières entre occurrences sont ici imposées : on ne compte pas de
paire reliant la fin d'une occurrence au début de la suivante.
Cette hypothèse explicite évite d'attribuer au calcul des paires inexistantes
dans l'exemple.

La taille de la séquence diminue, tandis que le vocabulaire gagne des entrées.
Les anciens symboles restent utiles pour des mots différents. Si l'on
rencontre « bac », la fusion apprise \((b,a)\) permet \([ba,c]\),
mais le symbole \code{baba} ne s'applique pas à toute la chaîne.
Apprendre BPE ne revient donc pas à enregistrer seulement une liste de mots.

\notion{L'ordre des fusions fait partie de l'apprentissage.}
Lorsqu'on ajoute une fusion, les comptages pertinents évoluent.
Des égalités de fréquence peuvent nécessiter une règle de départage.
Notre exemple possède des maxima simples ; il ne prétend pas prédire
les IDs ni toutes les décisions internes de la bibliothèque.
Le code du dépôt délègue ces mécanismes à \code{tokenizers}, au lieu
de réimplémenter les compteurs et leurs mises à jour.

\attention{La compression n'est pas un score de compréhension.}
Un symbole long très fréquent peut économiser des positions sans porter
un sens linguistique stable. Le Transformer apprendra ensuite à prédire
les unités obtenues. Changer de tokenizer change aussi l'unité de sa perte.

\exercice{Avec deux occurrences de « abac », compter les paires initiales.
Si l'on impose de fusionner \((a,b)\), quelle représentation obtient-on ?}
\corrige{Les paires \((a,b)\), \((b,a)\) et \((a,c)\) apparaissent chacune
deux fois : il y a égalité, donc aucune fusion unique n'est déductible des
seules fréquences. Avec le choix imposé, chaque occurrence devient
\([ab,a,c]\), soit trois tokens au lieu de quatre.}
